\subsubsection{Subproceso 2: Selección y Registro de Cursos}



\begin{enumerate}
    \item \textbf{Acceso autenticado a la Intranet DIRCE en turno habilitado.} \\
    \textit{Responsable: Estudiante Regular.} \\
    El estudiante ingresa al portal institucional con sus credenciales \texttt{@uni.pe} dentro de la fecha y franja horaria asignada, y el sistema verifica su estado de habilitación antes de mostrar el módulo de selección.

    \item \textbf{Visualización de la oferta académica personalizada.} \\
    \textit{Responsable: Sistema SIGA-ORCE / Intranet DIRCE.} \\
    La plataforma carga y despliega las asignaturas, secciones, horarios y cupos disponibles correspondientes al Plan de Estudios del estudiante, filtrando únicamente las que cumplen las condiciones de prerrequisito aprobadas.

    \item \textbf{Selección de asignaturas, secciones y horarios.} \\
    \textit{Responsable: Estudiante Regular.} \\
    El estudiante elige las asignaturas y secciones de su preferencia, priorizando las materias de arrastre del ciclo más bajo y los cursos en condición de riesgo académico, hasta completar la carga deseada.

    \item \textbf{Validación interactiva en tiempo real de restricciones.} \\
    \textit{Responsable: Sistema SIGA-ORCE / Intranet DIRCE.} \\
    El sistema comprueba automáticamente, por cada asignatura añadida, el cumplimiento del límite crediticio según el promedio ponderado y la ausencia de cruces de horario prohibidos, bloqueando o alertando al estudiante ante cualquier infracción.

    \item \textbf{Ajuste de la selección ante alertas del sistema.} \\
    \textit{Responsable: Estudiante Regular.} \\
    En caso de que el portal emita una alerta de restricción (cupo agotado, cruce inválido o exceso de créditos), el estudiante modifica su selección buscando secciones alternativas o replanteando la carga académica.

    \item \textbf{Confirmación y envío de la propuesta de matrícula.} \\
    \textit{Responsable: Estudiante Regular.} \\
    Una vez conformada la selección final, el estudiante ejecuta la acción de confirmación definitiva en el portal, formalizando el envío de su propuesta académica.

    \item \textbf{Registro de la Pre-Ficha y reserva de vacantes.} \\
    \textit{Responsable: Sistema SIGA-ORCE / Intranet DIRCE.} \\
    El sistema asienta la Pre-Ficha de Matrícula Registrada con el detalle de asignaturas, secciones y créditos, ejecuta el bloqueo transaccional de cupos para las secciones seleccionadas y pone la propuesta a disposición de la etapa de auditoría.
\end{enumerate}



